\vspace{-0.5em}
\section{Background}\label{sec:back}
Trust assessment in NNs requires a formalization that may represent ambiguity, inadequate evidence, and source reliability. Section~\ref{sec:backsl} summarizes the principles of Subjective Logic (SL) and the trust opinion representation, which encodes trust, distrust, and uncertainty. Section~\ref{sec:backslop} covers the main SL reasoning operators and their application in Subjective Trust Networks. Finally, we explain in Section~\ref{sec:backtwdt} how these concepts are applied to quantify dataset trustworthiness.
\vspace{-0.5em}
\subsection{Subjective Logic Fundamentals}\label{sec:backsl}
Reasoning about trust in AI systems requires handling partial, conflicting, or missing evidence. Classical probability theory models aleatoric uncertainty but cannot represent missing knowledge (ignorance) or contradictory information, which often arises from noisy or biased data.
Subjective Logic~\cite{josang2016subjective} extends Dempster–Shafer theory~\cite{dempster1967upper,shafer1976mathematical} to capture these conditions by representing trust as SL opinions rather than probabilities. A subjective opinion expresses beliefs and uncertainty as separate components, thereby distinguishing between uncertainty due to lack of knowledge (epistemic uncertainty) and uncertainty inherent to the environment itself (aleatoric uncertainty).

A subjective opinion denoted by \(\omega_X^A\) expresses the beliefs of an agent \(A\) (e.g., a sensor, a human, or an external observer of a process) about states of a variable $X$ which takes its values from a domain $\mathbb{X}$ (i.e., a state space). A special case of a subjective opinion is a subjective binomial opinion where $card(\mathbb{X})=2$. For a binary variable $X \in \mathbb{X} = \{x, \overline{x}\}$, a binomial opinion is expressed as a quadruple:
\vspace{-0.5em}
\begin{equation}
\omega_x = \omega_{X=x} = (b_x, d_x, u_x, a_x)
\end{equation}
satisfying $b_x + d_x + u_x = 1$, where $b_x$ denotes belief in $x$, $d_x$ disbelief in $x$ (belief in \(\overline{x}\)), $u_x$ the uncertainty mass, and $a_x$ the base rate (prior probability of $x$ in the absence of evidence). Its projected probability is defined as:
\vspace{-0.5em}
\begin{equation}
P(x) = b_x + a_x u_x.
\label{eq:projected_prob}
\end{equation}
This projection reduces the richer subjective opinion to an equivalent classical probability, enabling compatibility with standard probabilistic reasoning.

Binomial opinions can be derived from evidence through various quantification models. Let $r_x$ and $s_x$ denote the amount of positive and negative evidence, i.e.\ observations supporting respectively the truth and the falsehood of $x$. This work uses \emph{Baseline-Prior Quantification},
\begin{align}
b_x &= \frac{r_x}{W + r_x + s_x}, \quad
d_x = \frac{s_x}{W + r_x + s_x}, \quad
u_x = \frac{W}{W + r_x + s_x}
\label{eq:q1}
\end{align}
where the weight \(W>0\) guarantees residual uncertainty. Evidence-weighted and constant-uncertainty quantification~\cite{10.1007/978-3-032-19096-3_17} are alternatives that scale or fix the uncertainty mass instead.

A subjective opinion is meaningful only within a specific context or property under evaluation (e.g., accuracy, bias, or other trust-related aspects). In this work, trust is represented as a subjective binomial opinion $(t, d, u)$, where $t$ denotes trust (belief), $d$ distrust (disbelief), and $u$ uncertainty. These fundamentals define how trust is represented and interpreted as subjective opinions. To make them operational, Subjective Logic provides operators for combining, revising, and discounting opinions, which we introduce next.

\vspace{-0.5em}
\subsection{Subjective Logic Operators}\label{sec:backslop}

SL provides three families of reasoning operators for manipulating opinions~\cite{8009820,10706345,josang2016subjective}, all used in this work.

\begin{definition}[SL Operators]\label{def:fusion}
Let \(A\) be an agent reasoning about a variable \(X\).

\emph{Fusion} combines opinions held by two sources \(P\) and \(Q\) about the same proposition, \(\omega^A_{x} = \omega^P_{x} \oplus \omega^Q_{x}\); the appropriate variant (consensus, averaging, weighted, or cumulative) depends on whether the sources are correlated or independent.

\emph{Trust discounting} propagates an opinion along a referral chain: if \(A\) holds referral trust \(\omega^A_B\) in \(B\), who holds \(\omega^B_x\), then \(\omega^{[A;B]}_{x} = \omega^A_B \otimes \omega^B_{x}\), so that evidence from a partially reliable source is weakened and its uncertainty increased.

\emph{Inferential operators} generalize Bayesian reasoning through a conditional opinion \(\omega^A_{Y \mid X}\): \emph{deduction} derives \(\omega^A_Y\) from \(\omega^A_X\), and \emph{abduction} derives \(\omega^A_X\) from \(\omega^A_Y\).
\end{definition}
\noindent A Detail, summary and symbols of all operators appears in Appendices~\ref{app:operators} and \ref{app:operator}.

Building on these operators, \emph{Subjective Trust Networks}~\cite{4622580,10.1007/978-3-032-08887-1_4} model trust relationships as subjective opinions and propagated using trust discounting and fusion. 

\vspace{-0.5em}
\subsection{Trustworthiness in the dataset}\label{sec:backtwdt}
% \todo[inline]{Paper Trustworthiness in Dataset}

Focusing back on machine learning systems, the quality and structure of the training dataset are essential for determining the performance, robustness, and fairness of machine learning models. Common issues, such as sampling biases, mislabeled instances, or lack of diversity in the data can degrade learned representations and hinder generalization, thereby reducing the reliability of the model’s outputs~\cite{arp2021dosdontsmachinelearning}. A recent study on dataset quality shows that even small proportions of mislabeled samples can substantially shift model rankings; for instance, on CIFAR-10, VGG11 trained on clean data can outperform VGG19 once the fraction of erroneous labels reaches about 5\%, illustrating how sensitive benchmark conclusions are to label errors~\cite{northcutt2021pervasivelabelerrorstest}. Thus, evaluating the trustworthiness of training data is a critical step in assessing the trustworthiness of an AI system. 

Subjective Logic has been applied effectively to model dataset trustworthiness~\cite{10.1007/978-3-032-19096-3_17,trust-aidatatrust}, treating the dataset as a collection of samples, each composed of an input vector and a corresponding label. Trustworthiness can then be assessed at three granularities: the \emph{dataset level}, capturing global effects such as class imbalance or sampling bias; the \emph{instance level}, capturing mislabeled or corrupted data points; and the \emph{input feature level}, capturing fine-grained variation within an input vector, as when an adversarial patch renders only part of an image untrustworthy~\cite{Vargas2020} or when features originate from sources of differing reliability.
